% =============================================================================
%  PARTIE III — Documentation technique
% =============================================================================

\chapter{Architecture du projet}
\label{ch:architecture}

\section{Arborescence}

\begin{lstlisting}[caption={Arborescence du dépôt.},numbers=none]
Icare/
|-- pydmxIcare.py          # bas niveau : trame Enttec + port serie
|-- pydmxlight.py          # descripteurs de projecteurs (Par LED, Strobe)
|-- run_webapp.py          # point d'entree (python run_webapp.py)
|-- webapp/
|   |-- __init__.py        # fabrique create_app()
|   |-- config.py          # constantes, UserSettings, journalisation
|   |-- services.py        # singletons (settings, dmx)
|   |-- utils.py           # helpers partages (noms de fichiers surs)
|   |-- dmx_controller.py  # machine a etats + thread de sortie DMX
|   |-- live_player.py     # lecteur de sequence + tap-tempo
|   |-- dsl.py             # analyseur du langage de sequences
|   |-- schemas.py         # validation / migration des fichiers JSON
|   |-- pages.py           # routes HTML
|   |-- api_experiment.py  # API onglet Direct
|   |-- api_sequence.py    # API editeur de sequence
|   |-- api_live.py        # API lecture Live
|   |-- api_health.py      # /api/health
|   |-- api_palettes.py    # favoris + palettes nommees
|   |-- api_settings.py    # preferences
|   |-- templates/         # gabarits Jinja2
|   `-- static/            # app.css, app.js, vendor/three.min.js
|-- Sequences/             # shows d'exemple
|-- tests/                 # suite pytest
|-- docs/                  # cette documentation (LaTeX)
|-- pyproject.toml         # build, ruff, pytest
`-- build_exe.ps1          # paquetage PyInstaller
\end{lstlisting}

\section{La fabrique d'application}

\code{create\_app()} construit et configure l'application Flask. Elle choisit
le dossier des gabarits et des fichiers statiques selon qu'on tourne depuis les
sources ou depuis un exécutable empaqueté (\code{sys.\_MEIPASS}), puis enregistre
les \emph{blueprints} un par un :

\begin{lstlisting}[language=Python,caption={Enregistrement des blueprints (\file{webapp/\_\_init\_\_.py}).}]
from .pages import pages_bp
from .api_experiment import api_experiment_bp
...
app.register_blueprint(pages_bp)
app.register_blueprint(api_experiment_bp)
app.register_blueprint(api_sequence_bp)
app.register_blueprint(api_live_bp)
app.register_blueprint(api_health_bp)
app.register_blueprint(api_palettes_bp)
app.register_blueprint(api_settings_bp)
\end{lstlisting}

Chaque famille de routes vit dans son propre module : on retrouve donc
naturellement une route \api{/api/experiment/...} dans
\file{api\_experiment.py}, et ainsi de suite.

\section{Les singletons de service}

Pour éviter les imports circulaires et garantir une seule instance partagée,
\file{webapp/services.py} expose :

\begin{itemize}[leftmargin=1.4em]
  \item \code{settings} : les préférences chargées depuis le disque ;
  \item \code{dmx} : l'unique \code{DMXController} du processus ;
  \item \code{SEQ\_DIR} / \code{PALETTE\_DIR} : les chemins des dossiers de
        données.
\end{itemize}

\begin{note}
Il n'existe qu'\textbf{un} contrôleur DMX et \textbf{un} lecteur Live pour tout
le serveur. C'est voulu : un seul fil pilote la ligne DMX.
\end{note}

\section{Configuration}

\file{webapp/config.py} distingue deux choses :

\begin{itemize}[leftmargin=1.4em]
  \item les \textbf{constantes} : plan de canaux (\code{Channels}) et cadences
        de boucle (\code{Timings}) ;
  \item l'\textbf{état persistant} : la dataclass \code{UserSettings}.
\end{itemize}

\code{UserSettings} sait se charger, se sauvegarder et mettre à jour des champs.
Au chargement, seules les clés connues sont conservées, ce qui rend le fichier
tolérant aux évolutions de schéma. La journalisation écrit à la fois sur la
console et dans un fichier rotatif.

\begin{table}[H]
  \centering
  \begin{tabular}{@{}lll@{}}
    \toprule
    \textbf{Cadence} & \textbf{Valeur} & \textbf{Rôle} \\
    \midrule
    \code{STEADY\_LOOP} & 0,05 s & Sortie stable, environ 20 Hz. \\
    \code{RANDOM\_LOOP} & 0,025 s & Environ 40 Hz pour des fondus fluides. \\
    \code{EFFECT\_LOOP} & 0,025 s & Environ 40 Hz pour les effets. \\
    \code{IDLE\_LOOP} & 0,10 s & Attente quand aucun mode n'est actif. \\
    \bottomrule
  \end{tabular}
  \caption{Cadences de la boucle DMX (\code{Timings}).}
\end{table}

\section{Le front-end}

\file{webapp/static/app.js} est un unique script sans dépendance (hormis
Three.js, fourni en local dans \file{vendor/}). Au chargement, il lit
\code{document.body.dataset.page} et n'active que le module correspondant
(\code{bindExperiment}, \code{bindSequence} ou \code{bindLive}). Deux sondages
réguliers maintiennent l'interface synchronisée avec l'état serveur :

\begin{itemize}[leftmargin=1.4em]
  \item \api{GET /api/experiment/status} toutes les 200 ms (onglet Direct) ;
  \item \api{GET /api/live/status} toutes les 250 ms (onglet Live).
\end{itemize}

\chapter{Le contrôleur DMX}
\label{ch:controleur}

\code{DMXController} (\file{webapp/dmx\_controller.py}) est le cœur temps réel.

\section{Des modes exclusifs, un overlay}

Un seul mode peut être actif à la fois. Le strobe d'overlay, lui, est
indépendant et se superpose.

\begin{center}
\begin{tikzpicture}[
    font=\small,
    mode/.style={draw=icareblue,fill=icareblue!10,rounded corners,inner sep=7pt},
    arrow/.style={-Latex,draw=icodecmt},
  ]
  \node[mode] (idle) {inactif};
  \node[mode,right=1.6cm of idle] (steady) {steady};
  \node[mode,above=0.8cm of steady] (strobe) {strobe};
  \node[mode,below=0.8cm of steady] (random) {random};
  \node[mode,below=0.8cm of random] (effect) {effect};
  \node[mode,below=0.8cm of effect] (audio) {audio};
  \draw[arrow] (idle) -- (steady);
  \draw[arrow] (steady) -- (strobe);
  \draw[arrow] (strobe) -- (random);
  \draw[arrow] (random) -- (effect);
  \draw[arrow] (effect) -- (audio);
  \draw[arrow,dashed,bend left=35] (strobe.north) to node[above,draw=none,fill=none,font=\scriptsize]{switch} (idle.north);
\end{tikzpicture}
\end{center}
\legende{\code{\_switch\_mode()} réinitialise tous les drapeaux puis active le
mode voulu : les modes sont mutuellement exclusifs.}

\section{Le thread de sortie}

Le thread \code{dmx-out} boucle sans fin et exécute une étape selon le mode
actif :

\begin{lstlisting}[language=Python,caption={Boucle principale du contrôleur (simplifiée).}]
while not self._stop_event.is_set():
    try:
        if self.enabled:          self._send_with_overlay(); time.sleep(Timings.STEADY_LOOP)
        elif self.strobe_enabled: self._step_strobe()
        elif self.random_enabled: self._step_random()
        elif self.effect_enabled: self._step_effect()
        elif self.audio_enabled:  time.sleep(Timings.IDLE_LOOP)
        else:                     time.sleep(Timings.IDLE_LOOP)
    except Exception:
        logger.exception("DMX loop error; pausing 0.5 s")
        time.sleep(0.5)
\end{lstlisting}

Chaque \emph{frame} passe par \code{\_send\_with\_overlay()}, qui applique
d'abord l'overlay éventuel puis envoie l'image. Les écritures série sont
encapsulées dans des \code{try/except} : une erreur transitoire est journalisée
et la boucle repart après une courte pause, au lieu de tuer le thread.

\section{Points de robustesse}

\begin{itemize}[leftmargin=1.4em]
  \item Toutes les entrées sont \textbf{bornées} : les composantes RGB,
        l'amplitude, le pan et le tilt sont contraints à 0--255.
  \item Le canal de base est validé (0--509) et refusé sinon.
  \item \code{pydmxIcare.send\_data()} ne fait rien si le port est fermé : on
        peut lancer l'application sans matériel, sans plantage.
  \item En cas d'erreur d'écriture série, le port est refermé pour permettre une
        reconnexion propre.
\end{itemize}

\chapter{Le lecteur de séquence (LivePlayer)}
\label{ch:liveplayer}

\code{LivePlayer} (\file{webapp/live\_player.py}) parcourt les lignes d'une
séquence et traduit chaque cellule en couleurs.

\section{Chargement}

\code{load(raw)} normalise la séquence (voir chapitre~\ref{ch:schemas}) puis
remet l'index et l'état de lecture à zéro. Aucune lecture ne démarre tant que
\code{start()} n'est pas appelé.

\section{Transport}

\begin{table}[H]
  \centering
  \begin{tabular}{@{}lp{9.4cm}@{}}
    \toprule
    \textbf{Méthode} & \textbf{Effet} \\
    \midrule
    \code{start(bpm)} & Démarre la lecture (BPM optionnel). \\
    \code{pause()} & Bascule l'état \emph{pause}. \\
    \code{stop()} & Arrête et revient au début. \\
    \code{next()} / \code{prev()} & Avance ou recule d'une ligne. \\
    \code{toggle\_module(name)} & (Dé)sélectionne un module. \\
    \code{select\_all()} / \code{unselect\_all()} & Sélection globale. \\
    \bottomrule
  \end{tabular}
  \caption{API de transport du lecteur.}
\end{table}

\section{Boucle de lecture}

Le thread \code{live-player} applique une ligne puis attend la durée d'un temps :
\code{60 / BPM} secondes. À la fin de la séquence, la dernière image est
maintenue avant l'arrêt.

\begin{lstlisting}[language=Python,caption={Cœur de la boucle de lecture.}]
self.index = min(self.index + 1, total)
line = self.data["data"].get(str(self.index), {})
self._apply_step(line.get("function", []))
time.sleep(60.0 / max(1, self.bpm))
if self.index >= total:
    self.playing = False
\end{lstlisting}

\section{Sélection de modules et visualisation}

Si l'utilisateur a coché des modules, seuls ceux-là sont appliqués ; sinon tous
le sont. Pour la vue 3D, le lecteur retient la couleur du \emph{premier} module
ciblé et la pousse au contrôleur DMX, qui est mono-projecteur.

\section{Tap-tempo}

Le tempo mesuré vit dans \file{live\_player.py}. On conserve les horodatages
récents (fenêtre de 3 secondes), on calcule les intervalles, on en fait la
moyenne, et on borne le résultat entre 20 et 300 BPM :

\begin{lstlisting}[language=Python,caption={Calcul du tempo à partir des appuis.}]
intervals = [b - a for a, b in zip(self._taps, self._taps[1:], strict=False)]
avg = sum(intervals) / len(intervals)
bpm = int(round(60 / avg))
bpm = max(20, min(300, bpm))
\end{lstlisting}

\chapter{Le langage des séquences (DSL)}
\label{ch:dsl}

Chaque cellule d'une ligne contient une chaîne qui décrit l'action d'\textbf{un}
projecteur pour ce temps. \file{webapp/dsl.py} la transforme en objet typé.

\section{Grammaire}

\begin{table}[H]
  \centering
  \begin{tabular}{@{}ll@{}}
    \toprule
    \textbf{Jeton} & \textbf{Signification} \\
    \midrule
    \token{""} & Ne rien changer (garder l'état courant). \\
    \token{|} & Continuer l'effet multi-temps précédent. \\
    \token{uni\_n\_R\_G\_B} & Couleur fixe pendant \code{n} temps. \\
    \token{fond\_n\_R1\_G1\_B1\_R2\_G2\_B2} & Fondu linéaire de $c_1$ vers $c_2$. \\
    \token{str\_n\_R\_G\_B\_freq} & Strobe à \code{freq} Hz pendant \code{n} temps. \\
    \bottomrule
  \end{tabular}
  \caption{Jetons du langage de séquences.}
  \label{tab:dsl}
\end{table}

La fréquence du strobe est optionnelle : sans elle, la valeur par défaut est
60 Hz. Les longueurs \code{n} inférieures à 1 sont ramenées à 1, et toute
composante hors de 0--255 lève une \code{DSLError}.

\section{Structure du code}

Le module expose des dataclasses gelées :

\begin{itemize}[leftmargin=1.4em]
  \item \code{UniformEffect(length, color)} ;
  \item \code{FadeEffect(length, color\_from, color\_to)} ;
  \item \code{StrobeEffect(length, color, freq)} ;
  \item \code{ContinueEffect()} pour le jeton \token{|}.
\end{itemize}

La fonction \code{parse(token)} renvoie l'objet correspondant, ou \code{None}
pour une cellule vide. Trois \emph{encodeurs} (\code{encode\_uniform},
\code{encode\_fade}, \code{encode\_strobe}) font l'opération inverse, utilisés
notamment par la migration d'anciens fichiers.

\section{Le jeton de continuation}

Un effet de \code{n} temps s'écrit sur une première ligne, puis \code{n-1} fois
\token{|}. Le lecteur conserve, pour chaque module, un petit état
(mode, longueur, pas courant, couleurs, fréquence). À chaque \token{|}, il
avance d'un pas :

\begin{itemize}[leftmargin=1.4em]
  \item \textbf{fondu} : interpolation linéaire
        $c = c_1 + \frac{\text{pas}}{n-1}\,(c_2 - c_1)$ ;
  \item \textbf{strobe} : alternance couleur / noir à chaque pas ;
  \item \textbf{couleur fixe} : la couleur reste identique.
\end{itemize}

\begin{lstlisting}[caption={Exemple : un fondu de 4 temps du rouge au bleu.},numbers=none]
fond_4_255_0_0_0_0_255   |   |   |
\end{lstlisting}

\chapter{Schémas de fichiers et migration}
\label{ch:schemas}

Deux formats de séquence sont acceptés. \file{webapp/schemas.py} les unifie en
un seul : le format \emph{moderne}.

\section{Format moderne}

\begin{lstlisting}[language=json,caption={Séquence au format moderne.}]
{
  "Header": {
    "Structure": ["Intro_x4", "Intro_x4", "Couplet1_A_x2"],
    "Tempo": "60",
    "Modules": ["ParLED1", "ParLED2", "Strobe1"],
    "Line": ["1", "2", "3"]
  },
  "data": {
    "1": { "section": "Intro_x4", "function": ["", "uni_4_78_243_86", ""] },
    "2": { "section": "Intro_x4", "function": ["", "|", "str_3_132_214_107_63"] }
  }
}
\end{lstlisting}

\section{Format historique}

Les anciens fichiers (comme les premières versions de
\file{Doctor\_doctor.json}) ont une autre forme : une clé \code{headers} avec une
\code{structure}, et des lignes de quatre nombres
\code{[durée\_beats, R, G, B]}.

\begin{lstlisting}[language=json,caption={Ancien format (migré automatiquement).},numbers=none]
{
  "headers": { "structure": ["intro", "Couplet1", "Refrain1"] },
  "data": { "intro": [[2, 125, 20, 35], [2, 20, 53, 128]] }
}
\end{lstlisting}

\section{Normalisation}

\code{normalize\_sequence(raw)} renvoie toujours le format moderne :

\begin{itemize}[leftmargin=1.4em]
  \item les modules manquants sur une ligne sont complétés par des cellules
        vides, ou tronqués s'il y en a trop ;
  \item les clés de lignes sont converties en chaînes \code{"1"}, \code{"2"},
        \dots{} et triées numériquement ;
  \item en format historique, chaque durée est multipliée par deux (une durée de
        0,5 donne un temps, une durée de 2 en donne quatre) ; la première ligne
        reçoit un jeton \code{uni\_...}, les suivantes \token{|} ;
  \item un module factice \code{ParLED1} est utilisé, l'ancien format ne
        distinguant pas les projecteurs.
\end{itemize}

Si \code{raw} n'a ni \code{Header} ni \code{headers}, une \code{ValueError} est
levée : c'est ce qui alimente les réponses d'erreur 422 de l'API.

\begin{note}
À la sauvegarde, l'API valide en faisant un aller-retour de normalisation avant
d'écrire le fichier : un contenu invalide est donc refusé au lieu d'être
persisté.
\end{note}
